\chapter{Mortgage and Structured-Credit Strategies}\label{ch:s2:mortgage-and-structured-credit-strategies}

A homeowner can repay a mortgage whenever rates fall, so a mortgage bond is short an option that borrowers may not exercise well. A trader who models
the borrowers better than the market is paid for it, and one who models them worse pays. Gabaix, Krishnamurthy and Vigneron found that prepayment risk,
which washes out in the aggregate, is nonetheless priced, as if the marginal investor were a specialist rather than a diversified investor. On this
chapter's synthetic pool the market's prepayment model assumes mild burnout. An interest-only strip bought on the view that burnout is strong and
hedged with Treasuries gains 35.8\% of its price if the view is right. If borrowers never burn out, it loses 25.5\%. The build is
\texttt{firm.mbsrv}.

\section{Prepayment views}

\begin{definition}[Prepayment trade]\label{def:s2:mortgage-and-structured-credit-strategies:prepaytrade}
A \emph{prepayment trade}\index{prepayment trade} is a position in mortgage securities whose value depends on how fast borrowers repay (pass-throughs
of one coupon or pool type against another, interest-only or principal-only strips), taken because the trader's prepayment model differs from the one
in the market's prices, with the interest-rate risk hedged away.
\end{definition}

The market's model in \texttt{firm.mbsrv} is Book~6's (chapter~12): housing turnover on the PSA ramp, a refinancing S-curve in the incentive (the
borrowers' rate less the current mortgage rate) with a maximum CPR of 45\%, and burnout that decays refinancing by a factor $e^{-0.25 c}$, where $c$
is the cumulative positive incentive in percentage-point years. The pool pays its borrowers' 6.5\% to investors less 50 basis points, has 29 years
left, and is priced on Hull-White rate paths around a flat 4\% curve at an option-adjusted spread of 50 basis points. New mortgages cost about 5.8\%,
so the borrowers already have an incentive to refinance.

A trader believes these borrowers have passed up refinancing before and will keep passing it up: burnout of 1.0 instead of 0.25. If so, the pool
prepays more slowly than the market expects, and the interest-only strip, which is paid only while principal is outstanding, is worth more.

\section{Interest-only and principal-only trades}

The IO is worth 26.05 per 100 of pool face under the market's model. Its value rises with rates, since higher rates slow refinancing, so hedging it
means buying Treasuries: 0.51 ten-year Treasuries of face 100 per IO, the ratio that cancels the IO's rate risk as the market's model measures it
(\cref{lst:s2:mortgage-and-structured-credit-strategies:io}). The trade's gain is the IO's value under the model that turns out to be true, at the
same option-adjusted spread, less the price paid, plus the hedge.

\begin{center}
\small
\begin{tabular}{@{}lcc@{}}
\toprule
borrowers behave as & hedged gain, rates unchanged & rates 100 bp lower\\
\midrule
the market's model (burnout 0.25) & 0.0\% & 10.1\%\\
the view (burnout 1.0) & 35.8\% & 63.6\%\\
slower refinancing (maximum CPR 33\%) & 16.4\% & 31.1\%\\
never burning out (burnout 0) & $-25.5\%$ & $-36.0\%$\\
\bottomrule
\end{tabular}
\end{center}

The trade is a bet on borrowers, not on rates, but rates decide how much the bet matters
(\cref{fig:s2:mortgage-and-structured-credit-strategies:io}). When rates fall the incentive grows, and whether borrowers act on it is the whole
question. The view gains more and the wrong burnout assumption loses more. When rates rise the incentive disappears, and every model predicts the same
slow turnover. Even under the market's own model the hedged IO gains in large moves either way, because the hedge is sized for small ones.

\begin{omfigure}
\begin{tikzpicture}
\begin{axis}[width=11cm,height=5.6cm,xlabel={\small parallel shift of the curve (bp)},ylabel={\small hedged gain (\% of IO price)},
  tick label style={font=\footnotesize},xmin=-150,xmax=150,ymin=-45,ymax=90,grid=major,grid style={gray!25},table/col sep=comma,
  legend style={font=\scriptsize,at={(0.5,-0.3)},anchor=north},legend columns=2]
\addplot[black,thick] table[x=shift,y=market]{figdata/strategies-2/20-mortgage-and-structured-credit-strategies/io.csv};
\addplot[omThm,thick] table[x=shift,y=view]{figdata/strategies-2/20-mortgage-and-structured-credit-strategies/io.csv};
\addplot[gray,thick,dashed] table[x=shift,y=slow]{figdata/strategies-2/20-mortgage-and-structured-credit-strategies/io.csv};
\addplot[omDef,thick] table[x=shift,y=noburnout]{figdata/strategies-2/20-mortgage-and-structured-credit-strategies/io.csv};
\legend{market's model right,view right (burnout 1.0),slower refinancing,no burnout}
\end{axis}
\end{tikzpicture}

\omcaption{Gain of an interest-only strip bought at the market model's value and hedged with ten-year Treasuries, as a share of its price, against a
parallel shift of the curve, for four ways the borrowers might actually behave. Data: \texttt{s2\_mbsrv.io\_scenarios}.}
\label{fig:s2:mortgage-and-structured-credit-strategies:io}
\end{omfigure}

The principal-only strip is the mirror: it gains when prepayments are fast, so the trader who believes borrowers will refinance more than the market
expects buys the PO. The pair of strips lets a trader take either side of a prepayment view with the same pool.

\subsection*{The coupon stack}

The same view reprices every coupon. Pools whose borrowers pay more (higher coupons) have more incentive to refinance, so a model with stronger
burnout values them more. At the market's prices, the view's option-adjusted spread rises from 59 basis points on the 4.0\% coupon to 148 on the 7.5\%
(\cref{fig:s2:mortgage-and-structured-credit-strategies:oas}). Under a model without burnout, the same prices give OASs falling from 42 to $-326$
basis points. An OAS relative-value trade buys the coupon with the highest OAS under the trader's model and sells the lowest, hedged in duration. The
trade earns the OAS difference only if the model is right.

\begin{omfigure}
\begin{tikzpicture}
\begin{axis}[width=11cm,height=5.4cm,xlabel={\small pass-through coupon (\%)},ylabel={\small OAS at the market's price (bp)},
  tick label style={font=\footnotesize},xmin=4,xmax=7.5,ymin=-350,ymax=180,ytick={-300,-200,-100,0,100},grid=major,grid style={gray!25},table/col sep=comma,
  legend style={font=\scriptsize,at={(0.03,0.05)},anchor=south west}]
\addplot[black,thick] table[x=coupon,y=market]{figdata/strategies-2/20-mortgage-and-structured-credit-strategies/oas.csv};
\addplot[omThm,thick,mark=*,mark size=1.2pt] table[x=coupon,y=view]{figdata/strategies-2/20-mortgage-and-structured-credit-strategies/oas.csv};
\addplot[omDef,thick,mark=square*,mark size=1.2pt] table[x=coupon,y=noburnout]{figdata/strategies-2/20-mortgage-and-structured-credit-strategies/oas.csv};
\legend{market's model,view (burnout 1.0),no burnout}
\end{axis}
\end{tikzpicture}

\omcaption{The option-adjusted spread of each coupon at the price the market's model gives it (50 basis points), recomputed under the view's model and
under a model without burnout. Data: \texttt{s2\_mbsrv.stack}.}\label{fig:s2:mortgage-and-structured-credit-strategies:oas}
\end{omfigure}

Real OASs are not flat across coupons either. Boyarchenko, Fuster and Lucca found that MBS spreads show a pronounced smile across coupons, and
attributed the smile to prepayment risk that is not interest-rate risk, identified from the prices of stripped MBS. The time-series variation of
spreads, by contrast, was mostly explained by risks other than prepayment. A smile across coupons is what a market pricing borrowers' behaviour as a
risk, not a known function, would produce; it is also what a trader with the wrong model would see as mispricing.

\section{Tranche relative value}

\begin{definition}[Tranche relative value]\label{def:s2:mortgage-and-structured-credit-strategies:tranchrv}
\emph{Tranche relative value}\index{tranche relative value} is trading one tranche of a pooled credit structure against another, or against the
index or the underlying names, because the trader's model of default correlation or of the pool's losses differs from the one in the tranches'
prices.
\end{definition}

The large-pool Gaussian copula of Book~2 (chapter~24) prices tranches from the pool's default probability and one correlation. With a five-year
default probability of 6\% and 40\% recovery, the pool's expected loss is 3.6\% whatever the correlation. The tranches split it very differently as
correlation changes:

\begin{center}
\small
\begin{tabular}{@{}lccc@{}}
\toprule
expected loss, \% of tranche & correlation 0.15 & 0.30 & 0.45\\
\midrule
equity 0--3\% & 74.45 & 59.86 & 48.11\\
mezzanine 3--7\% & 25.32 & 24.34 & 22.11\\
mezzanine 7--10\% & 7.52 & 11.69 & 13.00\\
senior 10--15\% & 2.10 & 5.70 & 7.92\\
super senior 15--30\% & 0.15 & 1.23 & 2.77\\
\bottomrule
\end{tabular}
\end{center}

The senior tranche's expected loss is eighteen times larger at a correlation of 0.45 than at 0.15; the pool's is unchanged. Coval, Jurek and Stafford
named this fragility as one of two features that explain the rise and fall of structured finance: ratings that change drastically with modest errors
in the underlying risks, and exposure to systematic risk. In a companion paper they found that many structured products paid off like economic
catastrophe bonds, defaulting only in severe downturns, while offering far less compensation than alternatives with comparable payoffs.

A trader who believes the correlation is 0.15 while the market prices at 0.30 sells protection on the 7--10\% tranche and is paid the market's
expected loss of 11.7\% of the tranche over five years
(\cref{lst:s2:mortgage-and-structured-credit-strategies:tranche}). On 4,000 simulated pools of 125 names:

\begin{center}
\small
\begin{tabular}{@{}lccc@{}}
\toprule
true correlation & mean gain (\% of tranche) & chance of a loss & tranche wiped out\\
\midrule
0.15 (the view) & 3.1 & 11.3\% & 5.4\%\\
0.30 (the market) & $-0.4$ & 14.5\% & 9.4\%\\
0.45 & $-1.4$ & 15.0\% & 10.9\%\\
\bottomrule
\end{tabular}
\end{center}

Even when the view is right, the seller loses the whole tranche in 5.4\% of pools, about one in twenty. The small loss at the market's own correlation comes from the finite
pool of 125 names, which the large-pool formula ignores. A tranche view is a view on the tail, and a few years of history cannot confirm it.

\section{Strategy files}

\begin{strategyfile}{IO long on a slow-prepayment view}\label{strat:s2:mortgage-and-structured-credit-strategies:io}
\sfield{Who pays you, and why} Holders of prepayment risk who price borrowers with a generic model; specialists are the marginal investors.
\sfield{Instruments and venues} Interest-only strips of agency pools; Treasuries or swaps to hedge.
\sfield{Signal} A prepayment model by pool characteristics (loan size, seasoning, past incentive) against the market's speeds.
\sfield{Sizing and execution} Hedge the rate risk the market's model measures; size for the model being wrong.
\sfield{Costs} Wide bid-ask on strips; hedge rebalancing.
\sfield{How it dies} A refinancing wave the model underestimates; a policy change that makes refinancing easier.
\sfield{Horizon, capacity, infrastructure} Years; loan-level data and a prepayment model.
\sfield{Backtest honestly} Out-of-sample speeds by cohort; the model re-estimated only on data available at the time.
\sfield{Sources} Gabaix, Krishnamurthy and Vigneron (2007); this chapter: $+35.8\%$ if right, $-25.5\%$ with no burnout.
\end{strategyfile}

\begin{strategyfile}{OAS relative value across coupons}\label{strat:s2:mortgage-and-structured-credit-strategies:oas}
\sfield{Who pays you, and why} Investors who hold coupons for reasons other than their OAS: banks, the central bank, index trackers.
\sfield{Instruments and venues} TBAs and specified pools across the coupon stack; dollar rolls.
\sfield{Signal} Each coupon's OAS under the trader's model.
\sfield{Sizing and execution} Long high-OAS, short low-OAS coupons, duration-hedged.
\sfield{Costs} Roll financing; specified-pool payups.
\sfield{How it dies} The model's error, which is the OAS; supply from new issuance at one coupon.
\sfield{Horizon, capacity, infrastructure} Months.
\sfield{Backtest honestly} OAS from the model as it was, not refitted.
\sfield{Sources} Boyarchenko, Fuster and Lucca (2019); this chapter: an 88 basis point OAS difference under the view.
\end{strategyfile}

\begin{strategyfile}{Mezzanine tranche relative value}\label{strat:s2:mortgage-and-structured-credit-strategies:mezz}
\sfield{Who pays you, and why} Buyers or sellers of protection at one level of a structure for reasons other than its price: rating constraints,
hedging, regulatory capital.
\sfield{Instruments and venues} Index tranches; the index; CLO or CMBS tranches in cash form.
\sfield{Signal} Correlation or pool-loss views against the tranches' implied correlations.
\sfield{Sizing and execution} Delta-hedged with the index; size for the wipe-out.
\sfield{Costs} Wide bid-ask; model risk in the hedge ratio.
\sfield{How it dies} A systematic default wave; correlation that jumps.
\sfield{Horizon, capacity, infrastructure} Years; a copula and a pool-loss model.
\sfield{Backtest honestly} Stress the tail; history is too short to confirm a correlation view.
\sfield{Sources} Coval, Jurek and Stafford (2009); this chapter: $+3.1\%$ mean, wiped out in 5.4\% of pools.
\end{strategyfile}

\section{Tutorial: better than the market's model}\label{tut:s2:mortgage-and-structured-credit-strategies:tut}

\begin{tutorial}
\textbf{Goal.} Price a pool, its strips and its coupons with one prepayment model, revalue them with another, and sell protection on a tranche at the
wrong correlation. \textbf{End state:} the three tables and two figures.
\begin{tutsteps}
\item \textbf{The IO trade and the coupon stack}.
\omcode{code/firm/mbsrv/firm_mbsrv.py}{83}{112}{Buying the IO on a view, hedged; each coupon's OAS under the view.}\label{lst:s2:mortgage-and-structured-credit-strategies:io}
\item \textbf{The tranche}.
\omcode{code/firm/mbsrv/firm_mbsrv.py}{115}{126}{Selling protection at the market's correlation, paid on pools at the true one.}\label{lst:s2:mortgage-and-structured-credit-strategies:tranche}
\item \textbf{The four truths}.
\omcode{code/strategies-2/20-mortgage-and-structured-credit-strategies/python/s2_mbsrv.py}{32}{46}{The market's model, the view and two alternatives.}\label{lst:s2:mortgage-and-structured-credit-strategies:truths}
\item \textbf{Run} \texttt{io\_summary()}, \texttt{stack()}, \texttt{tranches()}, \texttt{mezz()} and \texttt{fig\_mbsrv.py}.
\end{tutsteps}
\textbf{What to change next.} Draw the true burnout at random and measure the trade's distribution; hedge the IO with the market's model and then
with the view's; delta-hedge the tranche with the index.
\end{tutorial}

\section{Build: mortgage and structured-credit relative value}\label{bld:s2:mortgage-and-structured-credit-strategies:mbsrv}

\begin{build}
\bfield{Purpose} Prepayment-view trades on Book 6's OAS model and Book 2's tranche model.
\bfield{Interface} \texttt{MbsConfig(\dots)}, \texttt{FlatCurve(rate)}, \texttt{strip\_paths(cfg, model, wac, oas, shift)},
\texttt{io\_trade(cfg, market, truth, side, shifts)}, \texttt{coupon\_stack(cfg, market, view, wacs)}, \texttt{tranche\_rv(p, rho\_market,
rho\_true, \dots)}.
\bfield{Rules} Prices from the market's model at one OAS; values under the true model at the same OAS; hedges sized with the market's model.
\bfield{Acceptance tests} \texttt{code/firm/mbsrv/tests/}: IO plus PO is the pass-through; no gain when the truth is the market's model; the same
model returns the market's OAS; equal correlations give equal expected losses.
\bfield{Stretch} Loan-level pools; dollar rolls; CLO waterfalls.
\end{build}

\begin{omsources}
\item X.~Gabaix, A.~Krishnamurthy and O.~Vigneron, ``Limits of arbitrage: theory and evidence from the mortgage-backed securities market'',
\emph{Journal of Finance} 62(2), 2007.
\item N.~Boyarchenko, A.~Fuster and D.~O.~Lucca, ``Understanding mortgage spreads'', \emph{Review of Financial Studies} 32(10), 2019.
\item J.~Coval, J.~Jurek and E.~Stafford, ``The economics of structured finance'', \emph{Journal of Economic Perspectives} 23(1), 2009.
\item J.~Coval, J.~Jurek and E.~Stafford, ``Economic catastrophe bonds'', \emph{American Economic Review} 99(3), 2009.
\end{omsources}

\section{Exercises}

\begin{exercise}[$\star$]\label{exo:s2:mortgage-and-structured-credit-strategies:1}
An IO bought at 26.05 per 100 of face gains 35.8\% of its price. What is the gain per 100 of face?
\end{exercise}

\begin{exercise}[$\star$]\label{exo:s2:mortgage-and-structured-credit-strategies:2}
Under the view, the 7.5\% coupon's OAS is 148 basis points and the 4.0\% coupon's 59; without burnout they are $-326$ and 42. What OAS difference
does a long 7.5\%, short 4.0\% trade earn in each case?
\end{exercise}

\begin{exercise}[$\star$]\label{exo:s2:mortgage-and-structured-credit-strategies:3}
The 7--10\% tranche of an index with notional 1 billion has an expected loss of 11.69\% of the tranche. What is the expected loss in money?
\end{exercise}

\begin{exercise}[$\star\star$]\label{exo:s2:mortgage-and-structured-credit-strategies:4}
Why does an IO lose value when rates fall, and why is hedging it a purchase of Treasuries?
\end{exercise}

\begin{exercise}[$\star\star$]\label{exo:s2:mortgage-and-structured-credit-strategies:5}
Why would prepayment risk, a wash in the aggregate, carry a price?
\end{exercise}

\begin{exercise}[$\star\star$]\label{exo:s2:mortgage-and-structured-credit-strategies:6}
Why is a senior tranche's rating fragile when the pool's expected loss is well estimated?
\end{exercise}

\begin{exercise}[$\star\star\star$]\label{exo:s2:mortgage-and-structured-credit-strategies:7}
\emph{Coding.} Rerun the IO trade with \texttt{MbsConfig(sigma=0.012)}. How do the price, the view's gain and the no-burnout loss change?
\end{exercise}

\begin{exercise}[$\star\star\star$]\label{exo:s2:mortgage-and-structured-credit-strategies:8}
\emph{Find the flaw.} ``The 7.5\% coupon's OAS is 148 basis points against 59 for the 4.0\%: it is 88 basis points cheap, and we lock that in.''
\end{exercise}

\section{Problem: Better Than the Market's Model}

\begin{problem}[{Weekend problem --- mortgage and structured credit}]\label{pb:s2:mortgage-and-structured-credit-strategies:1}
The chapter's synthetic pool and tranches and the public record.

\medskip\noindent\textbf{Part I --- The pool.}
\begin{enumerate}
\item Define a \omterm{def:s2:mortgage-and-structured-credit-strategies:prepaytrade}{prepayment trade}.
\item Describe the market's prepayment model and the pool.
\item What is burnout, and what does the view assume?
\item What did Gabaix, Krishnamurthy and Vigneron find?
\end{enumerate}

\medskip\noindent\textbf{Part II --- Strips and coupons.}
\begin{enumerate}[resume]
\item What is the IO worth, and how is it hedged?
\item Give the hedged gain under each of the four truths.
\item Why do rates decide how much the view matters?
\item How does the view change the coupon stack's OAS?
\end{enumerate}

\medskip\noindent\textbf{Part III --- Tranches.}
\begin{enumerate}[resume]
\item Define \omterm{def:s2:mortgage-and-structured-credit-strategies:tranchrv}{tranche relative value}.
\item How do tranche expected losses change with correlation?
\item What did Coval, Jurek and Stafford find?
\item Give the results of selling the 7--10\% tranche.
\end{enumerate}

\medskip\noindent\textbf{Part IV --- The verdict.}
\begin{enumerate}[resume]
\item State the \emph{named result}: the IO trade's return when the view is right and when the burnout assumption is wrong.
\item What does Boyarchenko, Fuster and Lucca's smile suggest about the market's model?
\item Why does the hedged IO gain in large moves even under the market's model?
\item How would you test a prepayment model honestly?
\item What would you size the IO trade for?
\item Why can history not confirm a tranche correlation view?
\item Which strategy file is closest to chapter~19's \omterm{def:s2:systematic-credit-and-bond-etf-arbitrage:factor}{credit factors}?
\item In one sentence: what does a mortgage trader sell?
\end{enumerate}
\end{problem}

\section{Interview questions}

\begin{interviewq}[$\star$ \iqroles{trader}]\label{iq:s2:mortgage-and-structured-credit-strategies:1}
What is negative convexity in a mortgage-backed security?
\end{interviewq}

\begin{interviewq}[$\star\star$ \iqroles{researcher}]\label{iq:s2:mortgage-and-structured-credit-strategies:2}
How would you build a prepayment model, and what variables would you use?
\end{interviewq}

\begin{interviewq}[$\star\star$ \iqroles{trader}]\label{iq:s2:mortgage-and-structured-credit-strategies:3}
You are long an IO. Rates fall 100 basis points. What happens, and what do you do?
\end{interviewq}

\begin{interviewq}[$\star\star$ \iqroles{risk}]\label{iq:s2:mortgage-and-structured-credit-strategies:4}
How would you measure the risk of a book of tranches beyond its deltas?
\end{interviewq}

\begin{interviewq}[$\star\star$ \iqroles{developer}]\label{iq:s2:mortgage-and-structured-credit-strategies:5}
An OAS calculation takes a minute per pool. How would you make it fast enough for a book of ten thousand pools?
\end{interviewq}

\begin{interviewq}[$\star\star\star$ \iqroles{researcher}]\label{iq:s2:mortgage-and-structured-credit-strategies:6}
In the large-pool Gaussian copula, show that the pool's expected loss does not depend on the correlation, and explain why a tranche's does.
\end{interviewq}
